\setgebiet{8}{Geschichte}{tief}
\chapter{Mathematikgeschichte}
\marginnote{\small\itshape\color{grau}\textbf{Sie brauchen:} Band 7, Kapitel 1--13. Die Rechnungen selbst brauchen nur Schulmathematik (Gebiet 0), Teilbarkeit (Gebiet 1) und für die Kapitel 9--11 Reihen und Gruppen (Gebiete 3 und 2).}[12pt]
\worum{Geschichte der Mathematik übt man nicht durch Auswendiglernen von Jahreszahlen, sondern indem man die alten Verfahren selbst ausführt: ägyptisch dividieren, babylonisch quadratische Gleichungen lösen, mit Archimedes $\pi$ einschließen, mit Cardano kubische Gleichungen knacken. Dabei sieht man, was die Alten konnten, was ihnen fehlte --- und welche Irrtümer sie machten, die man heute noch macht.}

\section{Was Sie hier üben}
\begin{center}
{\footnotesize\sffamily\setlength{\tabcolsep}{4pt}
\renewcommand{\arraystretch}{1.5}
\begin{tabular}{@{}lP{4.9cm}lcll@{}}
\toprule
\textbf{Nr.} & \textbf{Titel} & \textbf{Sorte} & \textbf{Grad} & \textbf{Zeit} & \textbf{Band/Kap.}\\
\midrule
7.1 & Ägyptische Stammbrüche & \typtypisch & \grad2 & 15 min & 7/1\\
7.2 & Mit und ohne Stellenwert & \typtypisch & \grad1 & 10 min & 7/1\\
7.3 & Ägyptisch dividieren & \typwerk & \grad1 & 10 min & 7/1\\
7.4 & Aufgabe 50 des Rhind-Papyrus & \typwerk & \grad1 & 10 min & 7/1\\
7.5 & Die Tafel YBC 7289 & \typwerk & \grad2 & 15 min & 7/1\\
7.6 & Eine babylonische quadratische Gleichung & \typwerk & \grad1 & 10 min & 7/1\\
7.7 & Brahmagupta und die Null & \typstolper & \grad1 & 10 min & 7/2\\
7.8 & Fangcheng & \typwerk & \grad2 & 15 min & 7/2\\
7.9 & Brahmaguptas Formel & \typwerk & \grad1 & 10 min & 7/2\\
7.10 & Wissen ist nicht Beweisen & \typstolper & \grad1 & 10 min & 7/3\\
7.11 & Wechselwegnahme & \typtypisch & \grad2 & 15 min & 7/3\\
7.12 & Inkommensurabel, nach Aristoteles & \typwerk & \grad2 & 15 min & 7/3\\
7.13 & Exhaustion & \typwerk & \grad3 & 20 min & 7/3\\
7.14 & Euklid IX.20 & \typwerk & \grad1 & 10 min & 7/4\\
7.15 & Archimedes und das 96-Eck & \typwerk & \grad2 & 20 min & 7/4\\
7.16 & Die Kugel auf der Waage & \typwerk & \grad3 & 20 min & 7/4\\
7.17 & Ptolemäus: $\operatorname{crd}36^{\circ}$ & \typwerk & \grad2 & 15 min & 7/5\\
7.18 & Diophant II.8 & \typwerk & \grad2 & 10 min & 7/5\\
7.19 & Die sechs Typen al-Chwarizmis & \typtypisch & \grad2 & 15 min & 7/6\\
7.20 & Quadrat und zehn Wurzeln & \typwerk & \grad1 & 10 min & 7/6\\
7.21 & Chayyams Kegelschnitte & \typwerk & \grad3 & 20 min & 7/6\\
7.22 & Bombelli und die Wurzel aus $-121$ & \typtypisch & \grad3 & 20 min & 7/7\\
7.23 & Fibonaccis Kaninchen & \typwerk & \grad1 & 10 min & 7/7\\
7.24 & Cardanos Formel & \typwerk & \grad2 & 20 min & 7/7\\
7.25 & Napiers Logarithmen & \typwerk & \grad2 & 15 min & 7/7\\
7.26 & Descartes' Normalenmethode & \typtypisch & \grad3 & 20 min & 7/8\\
7.27 & Fermats Adäqualität & \typwerk & \grad2 & 10 min & 7/8\\
7.28 & Das Teilungsproblem & \typwerk & \grad2 & 15 min & 7/8\\
7.29 & Die Leibniz-Reihe & \typtypisch & \grad2 & 15 min & 7/9\\
7.30 & Was Euler durfte & \typstolper & \grad2 & 15 min & 7/9\\
7.31 & Newtons Binomialreihe & \typwerk & \grad2 & 15 min & 7/9\\
7.32 & Leibniz' Produktregel & \typwerk & \grad1 & 10 min & 7/9\\
7.33 & Das Basler Problem & \typwerk & \grad3 & 20 min & 7/9\\
7.34 & Ein Dedekind-Schnitt & \typtypisch & \grad2 & 15 min & 7/10\\
7.35 & Die Weierstraß-Funktion & \typwerk & \grad2 & 15 min & 7/10\\
7.36 & Cantors Diagonale & \typwerk & \grad2 & 15 min & 7/10\\
7.37 & Lagranges Resolvente & \typtypisch & \grad2 & 15 min & 7/11\\
7.38 & Die Galoisgruppe von $x^{3} - 2$ & \typwerk & \grad3 & 20 min & 7/11\\
7.39 & $2 + 2 = 4$ nach Peano & \typtypisch & \grad2 & 15 min & 7/12\\
7.40 & Die Menge aller Mengen & \typstolper & \grad2 & 15 min & 7/13\\
\bottomrule
\end{tabular}}
\end{center}

\section{Woran Sie merken, dass Sie es können}
\begin{itemize}[leftmargin=1.2em,itemsep=2pt]
\item Sie können ein altes Verfahren ausführen \emph{und} sagen, welcher moderne Satz dahintersteht --- und was den Alten zum Beweis fehlte.
\item Sie unterscheiden, was eine Kultur \emph{wusste} (Rechenregeln, Tabellen) von dem, was sie \emph{bewies}.
\item Sie erkennen, warum $\sqrt{-121}$, $\frac00$ und $1 - 1 + 1 - \dots$ jahrhundertelang Streitfragen waren, und wo die heutige Antwort steht.
\item Sie sehen in jeder modernen Definition ($\varepsilon$-$\delta$, Dedekind-Schnitt, Peano-Axiome) die Frage, die sie beantwortet.
\end{itemize}

\section{Aufgaben}

%%%%%%%%%%% Kapitel 1
\stupser{tief}{Gieriger Algorithmus: Immer den größten Stammbruch abziehen, der noch hineinpasst. $\frac4{13}$: $\frac14$ passt, Rest $\frac3{52}$; dann $\frac1{18}$; Rest $\frac1{468}$.}{22pt}
\begin{karte}{tief}{7.1}{\typtypisch}{\meta{2}{15 min}{Band 7, Kap.\ 1}}
\kartentitel{Ägyptische Stammbrüche}
\worum{Die Ägypter schrieben jeden Bruch als Summe verschiedener Stammbrüche $\frac1n$. Der Rhind-Papyrus beginnt mit einer Tabelle für $\frac2n$. Die Aufgabe übt das Verfahren --- und zeigt, dass die ägyptischen Zerlegungen besser sind als die naheliegenden.}
\methoden{Stammbruchzerlegung, Gieriger Algorithmus, Bruchrechnung}
\begin{enumerate}[label=(\alph*),leftmargin=1.6em,itemsep=2pt]
\item Prüfen Sie die Rhind-Zerlegungen $\frac27 = \frac14 + \frac1{28}$ und $\frac2{13} = \frac18 + \frac1{52} + \frac1{104}$.
\item Zerlegen Sie $\frac4{13}$ mit dem gierigen Verfahren (Fibonacci 1202): jeweils den größten passenden Stammbruch abziehen.
\item Zerlegen Sie $\frac2{13}$ gierig und vergleichen Sie mit der Rhind-Tabelle. Welche Zerlegung ist \enquote{besser}, und in welchem Sinn?
\item Warum endet das gierige Verfahren immer? (Betrachten Sie den Zähler des Rests.)
\end{enumerate}
\end{karte}

\stupser{tief}{Römisch multipliziert man nicht --- man verdoppelt wie die Ägypter: $\mathrm{XLVII}$, $\mathrm{XCIV}$, $\mathrm{CLXXXVIII}$, $\mathrm{CCCLXXVI}$, und $12 = 8 + 4$.}{22pt}
\begin{karte}{tief}{7.2}{\typtypisch}{\meta{1}{10 min}{Band 7, Kap.\ 1}}
\kartentitel{Mit und ohne Stellenwert}
\worum{Drei Zahlsysteme, drei Rechenweisen: additiv (römisch), sexagesimal mit Stellenwert (babylonisch), dezimal mit Null (indisch). Die Aufgabe macht fühlbar, was der Stellenwert leistet.}
\methoden{Zahlsysteme, Sexagesimalsystem, Verdopplungsverfahren}
\begin{enumerate}[label=(\alph*),leftmargin=1.6em,itemsep=2pt]
\item Berechnen Sie $\mathrm{XLVII}\cdot\mathrm{XII}$, ohne in Dezimalzahlen zu übersetzen --- durch Verdoppeln und Addieren römischer Zahlen.
\item Berechnen Sie babylonisch $1;30\cdot2;15$ (Semikolon trennt Ganze und Sechzigstel) und prüfen Sie dezimal.
\item Schreiben Sie $\frac13$ und $\frac17$ sexagesimal. Warum war die Sechzig eine gute Wahl --- und wo versagt sie?
\item Die Babylonier hatten lange kein Zeichen für die Null und keines für das Komma. Was bedeutet die Zahl \enquote{$1, 30$} dann? Nennen Sie drei Möglichkeiten.
\end{enumerate}
\end{karte}

\stupser{gruen}{Verdopplungstabelle für $31$: $31, 62, 124, 248, 496, 992$. Welche davon summieren sich zu höchstens $1876$?}{22pt}
\begin{karte}{gruen}{7.3}{\typwerk}{\meta{1}{10 min}{Band 7, Kap.\ 1, Aufgabe 1.1}}
\kartentitel{Ägyptisch dividieren}
\worum{Division als rückwärts gelesene Multiplikation, und Multiplikation als Verdoppeln --- das ägyptische Verfahren ist das binäre System, zweitausend Jahre vor Leibniz.}
\methoden{Verdopplungsverfahren, Binärdarstellung, Division mit Rest}
Berechnen Sie $1876 : 31$ nach ägyptischer Art (durch Verdoppeln). Was hat das mit der Binärdarstellung von $60$ zu tun?
\end{karte}

\stupser{gruen}{$(d - \frac d9)^{2} = (\frac89d)^{2} = \frac{64}{81}d^{2}$; vergleichen mit $\frac\pi4d^{2}$.}{22pt}
\begin{karte}{gruen}{7.4}{\typwerk}{\meta{1}{10 min}{Band 7, Kap.\ 1, Aufgabe 1.2}}
\kartentitel{Aufgabe 50 des Rhind-Papyrus}
\worum{Ein Rezept ohne Begründung, $1650$ v.\,Chr.: Vom Durchmesser ein Neuntel weg, den Rest quadrieren. Welches $\pi$ steckt darin?}
\methoden{Kreisfläche, Näherung bewerten, Rezept entschlüsseln}
Ein rundes Feld hat den Durchmesser $9$. Die Anweisung: Nimm $\frac19$ vom Durchmesser weg; der Rest ist $8$. Multipliziere $8$ mit $8$; es ergibt $64$. Welchen Wert für $\pi$ benutzt das Verfahren, und wie genau ist es? Wie könnte man auf $\frac89$ gekommen sein?
\end{karte}

\stupser{gruen}{$1;24,51,10 = 1 + \frac{24}{60} + \frac{51}{3600} + \frac{10}{216\,000}$. Und $42;25,35$ ist das $30$-fache davon.}{22pt}
\begin{karte}{gruen}{7.5}{\typwerk}{\meta{2}{15 min}{Band 7, Kap.\ 1, Aufgabe 1.3}}
\kartentitel{Die Tafel YBC 7289}
\worum{Eine Tontafel, um $1700$ v.\,Chr.: ein Quadrat, seine Diagonalen, drei Zahlen. Sie enthält $\sqrt2$ auf sechs Dezimalstellen --- und verrät, dass die Babylonier den Satz des Pythagoras kannten, tausend Jahre vor Pythagoras.}
\methoden{Sexagesimalsystem, Näherung bewerten, Satz des Pythagoras}
Die Tafel zeigt ein Quadrat mit den Diagonalen und den Zahlen $1;24,51,10$ und $42;25,35$. Deuten Sie die Tafel und bestimmen Sie die Genauigkeit. Mit welchem Verfahren könnten die Babylonier die Wurzel berechnet haben?
\end{karte}

\stupser{gruen}{Quadratische Ergänzung in Worten: Hälfte des Koeffizienten, quadrieren, addieren, Wurzel, Hälfte wieder dazu.}{22pt}
\begin{karte}{gruen}{7.6}{\typwerk}{\meta{1}{10 min}{Band 7, Kap.\ 1, Aufgabe 1.4}}
\kartentitel{Eine babylonische quadratische Gleichung}
\worum{BM 13901, Nr.\ 2: \enquote{Ich habe die Fläche eines Quadrats von seiner Seite subtrahiert, und es ist $870$.} Die Lösung ist die quadratische Ergänzung --- als Rechenrezept, ohne Formel und ohne Symbole.}
\methoden{Quadratische Ergänzung, Rezept entschlüsseln, Quadratische Gleichung}
Lösen Sie $x^{2} - x = 870$ mit dem babylonischen Verfahren (Schritt für Schritt in Worten, dann modern). Warum findet das Verfahren nur eine Lösung?
\end{karte}

%%%%%%%%%%% Kapitel 2
\stupser{akzent}{Brahmagupta (628): $a - a = 0$, $0\cdot a = 0$, $0 : 0 = 0$. Das erste und zweite sind richtig, das dritte nicht --- und $a : 0$ ließ er offen.}{22pt}
\begin{karte}{akzent}{7.7}{\typstolper}{\meta{1}{10 min}{Band 7, Kap.\ 2}}
\kartentitel{Brahmagupta und die Null}
\worum{Brahmagupta war der Erste, der mit der Null als Zahl rechnete --- und der Erste, der sich bei $\frac00$ irrte. Bhaskara setzte $500$ Jahre später $\frac a0 = \infty$. Beides ist falsch, und die Gründe sind lehrreich.}
\methoden{Division durch null vermeiden, Gegenbeispiel, Grenzwerte}
\begin{enumerate}[label=(\alph*),leftmargin=1.6em,itemsep=2pt]
\item Brahmagupta setzt $\frac00 = 0$. Zeigen Sie mit der Definition der Division ($\frac ab = c\iff a = bc$), dass $\frac00$ jeden Wert haben könnte.
\item Bhaskara setzt $\frac a0 = \infty$ für $a\neq0$. Was ist daran richtig gemeint, und warum ist es als Gleichung unhaltbar? (Betrachten Sie $\frac1x$ für $x\to0^{+}$ und $x\to0^{-}$.)
\item Jemand \enquote{beweist} $1 = 2$: Aus $a = b$ folgt $a^{2} = ab$, also $a^{2} - b^{2} = ab - b^{2}$, also $(a + b)(a - b) = b(a - b)$, also $a + b = b$, also $2b = b$. Wo steckt die Null?
\item Was ist $\lim_{x\to0}\frac{\sin x}x$, $\lim_{x\to0}\frac{x^{2}}x$, $\lim_{x\to0}\frac x{x^{2}}$? Was lehrt das über \enquote{$\frac00$}?
\end{enumerate}
\end{karte}

\stupser{gruen}{Spaltenweise anordnen, dann die rechte Spalte mit $3$ multiplizieren und die linke so oft abziehen, bis oben eine Null steht --- Gauß auf dem Rechenbrett.}{22pt}
\begin{karte}{gruen}{7.8}{\typwerk}{\meta{2}{15 min}{Band 7, Kap.\ 2, Aufgabe 2.1}}
\kartentitel{Fangcheng}
\worum{Die \emph{Neun Kapitel} lösen lineare Gleichungssysteme mit Rechenstäbchen in rechteckiger Anordnung --- der Gauß-Algorithmus, etwa $200$ v.\,Chr. Die Garbenaufgabe ist dieselbe wie Aufgabe 2.2; hier wird sie so gerechnet wie damals.}
\methoden{Gauß-Algorithmus, Rechenbrett, Lineare Gleichungssysteme}
$3$ Garben guter, $2$ mittlerer, $1$ schlechter Ernte ergeben $39$ Dou; $2$, $3$, $1$ ergeben $34$; $1$, $2$, $3$ ergeben $26$. Lösen Sie mit dem Verfahren \emph{fangcheng}: spaltenweise Anordnung, Spalten vervielfachen und voneinander abziehen. Wo tauchen negative Zahlen auf, und wie gingen die Chinesen damit um?
\end{karte}

\stupser{gruen}{$s = 9$: $A = \sqrt{6\cdot5\cdot4\cdot3} = \sqrt{360}$. Mit $d = 0$ wird aus dem Viereck ein Dreieck --- und aus Brahmagupta Heron.}{22pt}
\begin{karte}{gruen}{7.9}{\typwerk}{\meta{1}{10 min}{Band 7, Kap.\ 2, Aufgabe 2.2}}
\kartentitel{Brahmaguptas Formel}
\worum{Für Sehnenvierecke gilt $A = \sqrt{(s - a)(s - b)(s - c)(s - d)}$ --- die Heronsche Formel mit einer Seite mehr. Dass der Umkreis nötig ist, sieht man am Gegenbeispiel.}
\methoden{Heronsche Formel, Sehnenviereck, Spezialfall prüfen}
Berechnen Sie die Fläche eines Sehnenvierecks mit den Seiten $3, 4, 5, 6$ und prüfen Sie den Zusammenhang mit der Heron-Formel. Warum gilt die Formel nicht für beliebige Vierecke? (Ein Quadrat mit Seite $1$ und eine Raute mit Seite $1$ und Winkel $30^{\circ}$ haben dieselben Seiten.)
\end{karte}

%%%%%%%%%%% Kapitel 3
\stupser{akzent}{Plimpton 322 ($1800$ v.\,Chr.) listet pythagoreische Tripel wie $(119, 120, 169)$. Eine Tabelle ist kein Beweis --- aber auch nicht nichts.}{22pt}
\begin{karte}{akzent}{7.10}{\typstolper}{\meta{1}{10 min}{Band 7, Kap.\ 3}}
\kartentitel{Wissen ist nicht Beweisen}
\worum{\enquote{Pythagoras hat den Satz des Pythagoras entdeckt} --- nein: Babylonier, Inder und Chinesen kannten ihn vorher. \enquote{Also hat er ihn nicht bewiesen} --- auch nicht zwingend. Die Aufgabe übt, drei Dinge auseinanderzuhalten: eine Regel kennen, sie an Beispielen prüfen, sie beweisen.}
\methoden{Wissen und Beweis unterscheiden, Quellenkritik, Satz des Pythagoras}
\begin{enumerate}[label=(\alph*),leftmargin=1.6em,itemsep=2pt]
\item Prüfen Sie das Tripel $(119, 120, 169)$ aus Plimpton 322. Zeigt die Tafel, dass die Babylonier den Satz \emph{bewiesen} hatten?
\item Das chinesische \emph{Zhoubi Suanjing} zeigt eine Figur: ein Quadrat mit Seite $a + b$, darin vier rechtwinklige Dreiecke und ein schiefes Quadrat. Rekonstruieren Sie daraus einen Beweis. Ist die Figur ein Beweis?
\item Was unterscheidet Euklids Beweis (I.47) von beiden? Nennen Sie zwei Merkmale.
\item Formulieren Sie in einem Satz, was \enquote{die Griechen erfanden den Beweis} bedeutet und was nicht.
\end{enumerate}
\end{karte}

\stupser{tief}{Wechselwegnahme ist der euklidische Algorithmus mit Strecken: Die kleinere von der größeren abtragen, so oft es geht; mit dem Rest weiter. Bricht es ab, sind die Strecken kommensurabel.}{22pt}
\begin{karte}{tief}{7.11}{\typtypisch}{\meta{2}{15 min}{Band 7, Kap.\ 3}}
\kartentitel{Wechselwegnahme}
\worum{Die Griechen verglichen Strecken, indem sie die kleinere von der größeren wegnahmen --- \emph{anthyphairesis}. Endet das Verfahren, haben die Strecken ein gemeinsames Maß. Beim goldenen Schnitt endet es nie, und das sieht man sofort.}
\methoden{Euklidischer Algorithmus, Kettenbruch, Inkommensurabilität}
\begin{enumerate}[label=(\alph*),leftmargin=1.6em,itemsep=2pt]
\item Führen Sie die Wechselwegnahme für die Strecken $17$ und $5$ aus. Welches gemeinsame Maß ergibt sich, und welcher Kettenbruch?
\item Im goldenen Rechteck (Seiten $a > b$ mit $a : b = b : (a - b)$) nimmt man $b$ einmal von $a$ weg und erhält ein Rechteck mit Seiten $b$ und $a - b$. Was ist das für ein Rechteck? Warum endet die Wechselwegnahme nie?
\item Folgern Sie: Der goldene Schnitt ist irrational, und sein Kettenbruch ist $[1; 1, 1, 1, \dots]$.
\item Führen Sie dasselbe für Seite und Diagonale des Quadrats durch (Tipp: Diagonale minus Seite, dann im kleineren Quadrat). Welcher Kettenbruch?
\end{enumerate}
\end{karte}

\stupser{gruen}{Ist $\frac ds = \frac mn$ gekürzt, so $m^{2} = 2n^{2}$: $m$ gerade, also $m^{2}$ durch $4$ teilbar, also $n^{2}$ gerade, also $n$ gerade --- beide gerade, Widerspruch.}{22pt}
\begin{karte}{gruen}{7.12}{\typwerk}{\meta{2}{15 min}{Band 7, Kap.\ 3, Aufgabe 3.1}}
\kartentitel{Inkommensurabel, nach Aristoteles}
\worum{Der älteste überlieferte Widerspruchsbeweis: Aristoteles nennt ihn als Musterbeispiel und sagt, er führe darauf, dass eine Zahl zugleich gerade und ungerade sei. Hier in genau dieser Form.}
\methoden{Widerspruchsbeweis, Gerade und ungerade, Inkommensurabilität}
Führen Sie den antiken Beweis, dass Quadratseite und Diagonale inkommensurabel sind, in der Form, die Aristoteles überliefert. Warum ist das für die Pythagoreer eine Krise gewesen?
\end{karte}

\stupser{gruen}{Angenommen $A_1 : A_2 > d_1^{2} : d_2^{2}$. Dann gibt es eine Fläche $S < A_2$ mit $A_1 : S = d_1^{2} : d_2^{2}$; ein einbeschriebenes Vieleck im zweiten Kreis mit Fläche $> S$ führt zum Widerspruch.}{22pt}
\begin{karte}{gruen}{7.13}{\typwerk}{\meta{3}{20 min}{Band 7, Kap.\ 3, Aufgabe 3.2}}
\kartentitel{Exhaustion}
\worum{Eudoxos' Methode, mit der die Griechen Grenzwerte vermieden: Man zeigt, dass weder \enquote{größer} noch \enquote{kleiner} möglich ist. Euklid XII.2 ist das Musterbeispiel --- Kreisflächen verhalten sich wie die Quadrate der Durchmesser.}
\methoden{Exhaustionsmethode, Doppelter Widerspruch, Grenzwerte}
Beweisen Sie mit der Exhaustionsmethode, dass $A_1 : A_2 = d_1^{2} : d_2^{2}$ für zwei Kreise. Wo steckt der moderne Grenzwert, und wie umgeht Euklid ihn?
\end{karte}

%%%%%%%%%%% Kapitel 4
\stupser{gruen}{$N = p_1\cdots p_n + 1$ hat einen Primteiler; keiner der $p_i$ teilt $N$ (Rest $1$). Euklid sagt nicht \enquote{unendlich}, sondern \enquote{mehr als jede vorgelegte Anzahl}.}{22pt}
\begin{karte}{gruen}{7.14}{\typwerk}{\meta{1}{10 min}{Band 7, Kap.\ 4, Aufgabe 4.1}}
\kartentitel{Euklid IX.20}
\worum{Der berühmteste Beweis der Mathematik --- und eine Lektion in Formulierung: Euklid vermeidet das aktual Unendliche, und sein Beweis ist kein Widerspruchsbeweis.}
\methoden{Euklids Primzahlbeweis, Konstruktiver Beweis, Aktual und potentiell unendlich}
Beweisen Sie mit Euklid, dass es mehr Primzahlen gibt als jede vorgelegte Anzahl. Ist $N = p_1\cdots p_n + 1$ selbst immer prim? Warum ist Euklids Beweis konstruktiv, obwohl er meist als Widerspruchsbeweis erzählt wird?
\end{karte}

\stupser{gruen}{Rekursion für die halben Umfänge $a_n$ (ein-) und $b_n$ (umbeschrieben): $b_{2n} = \frac{2a_nb_n}{a_n + b_n}$, $a_{2n} = \sqrt{a_nb_{2n}}$; Start $a_6 = 3$, $b_6 = 2\sqrt3$.}{22pt}
\begin{karte}{gruen}{7.15}{\typwerk}{\meta{2}{20 min}{Band 7, Kap.\ 4, Aufgabe 4.2}}
\kartentitel{Archimedes und das 96-Eck}
\worum{$3\frac{10}{71} < \pi < 3\frac17$: Archimedes' Einschließung, um $250$ v.\,Chr., durch viermaliges Verdoppeln vom Sechseck zum $96$-Eck --- ohne Dezimalbrüche, mit Wurzelabschätzungen von Hand.}
\methoden{Einschließung, Rekursion, Mittelwerte}
Rechnen Sie Archimedes' $\pi$-Einschließung nach: vom Sechseck durch Verdoppeln zum $96$-Eck. Vergleichen Sie mit seinen Schranken $3\frac{10}{71}$ und $3\frac17$.
\end{karte}

\stupser{gruen}{Scheiben der Kugel und des Kegels bei $x$, an den Hebelarm $2r$ gehängt, halten den Zylinderscheiben bei $x$ das Gleichgewicht: $2r(\pi\rho^{2} + \pi x^{2}) = x\cdot\pi(2r)^{2}$ mit $\rho^{2} = 2rx - x^{2}$.}{22pt}
\begin{karte}{gruen}{7.16}{\typwerk}{\meta{3}{20 min}{Band 7, Kap.\ 4, Aufgabe 4.3}}
\kartentitel{Die Kugel auf der Waage}
\worum{Archimedes' \emph{Methode}, $1906$ wiederentdeckt: Er wog Kugel und Kegel gegen einen Zylinder --- Scheibe für Scheibe --- und fand so $V = \frac43\pi r^{3}$, bevor er es streng bewies. Integration als Mechanik.}
\methoden{Hebelgesetz, Scheibenmethode, Heuristik und Beweis}
Vollziehen Sie Archimedes' mechanische Bestimmung des Kugelvolumens nach. Warum nannte er das Verfahren selbst keinen Beweis?
\end{karte}

%%%%%%%%%%% Kapitel 5
\stupser{gruen}{Zehneckseite $=$ größerer Abschnitt des goldenen Schnitts von $R$: $x = R\frac{\sqrt5 - 1}2$. Also $\sin18^{\circ} = \frac{\sqrt5 - 1}4$.}{22pt}
\begin{karte}{gruen}{7.17}{\typwerk}{\meta{2}{15 min}{Band 7, Kap.\ 5, Aufgabe 5.1}}
\kartentitel{Ptolemäus: $\operatorname{crd}36^{\circ}$}
\worum{Ptolemäus' Sehnentafel ($150$ n.\,Chr.) ist die erste trigonometrische Tabelle. Ihr Grundstein: die Sehne von $36^{\circ}$ aus dem regelmäßigen Zehneck und dem goldenen Schnitt.}
\methoden{Sehnenrechnung, Goldener Schnitt, Trigonometrie}
Berechnen Sie mit Ptolemäus' Methoden $\operatorname{crd}36^{\circ}$ (die Sehne zum Zentriwinkel $36^{\circ}$ im Kreis mit Radius $R$) und daraus $\sin18^{\circ}$. Wie kam Ptolemäus von dort zu $\operatorname{crd}1^{\circ}$?
\end{karte}

\stupser{gruen}{Ansatz $y = 2t - 4$ (die $4$ ist $\sqrt{16}$, die $2$ frei gewählt): $t^{2} + (2t - 4)^{2} = 16$ wird linear in $t$, weil die $16$ wegfällt.}{22pt}
\begin{karte}{gruen}{7.18}{\typwerk}{\meta{2}{10 min}{Band 7, Kap.\ 5, Aufgabe 5.2}}
\kartentitel{Diophant II.8}
\worum{\enquote{Ein gegebenes Quadrat in zwei Quadrate zu zerlegen} --- die Aufgabe, an deren Rand Fermat seine Vermutung notierte. Diophants Trick macht aus einer quadratischen Gleichung eine lineare.}
\methoden{Rationale Parametrisierung, Diophantische Gleichung, Geschickter Ansatz}
Lösen Sie Diophants Aufgabe II.8 mit seiner Methode für $16 = x^{2} + y^{2}$. Was liefert der Ansatz $y = kt - 4$ allgemein, und was hat das mit dem Einheitskreis zu tun?
\end{karte}

%%%%%%%%%%% Kapitel 6
\stupser{tief}{Al-Chwarizmi kennt keine negativen Zahlen und keine Null als Koeffizient: Deshalb sechs Typen statt einer Gleichung $ax^{2} + bx + c = 0$.}{22pt}
\begin{karte}{tief}{7.19}{\typtypisch}{\meta{2}{15 min}{Band 7, Kap.\ 6}}
\kartentitel{Die sechs Typen al-Chwarizmis}
\worum{Das \emph{Kitab al-dschabr} ($820$) teilt quadratische Gleichungen in sechs Typen ein, alle mit positiven Koeffizienten. Die Aufgabe übt die Einteilung --- und zeigt, wo zwei Lösungen auftauchen und wo eine verschwindet.}
\methoden{Quadratische Gleichung, Fallunterscheidung, Historische Normalform}
Die sechs Typen: (1) $ax^{2} = bx$, (2) $ax^{2} = c$, (3) $bx = c$, (4) $x^{2} + bx = c$, (5) $x^{2} + c = bx$, (6) $x^{2} = bx + c$.
\begin{enumerate}[label=(\alph*),leftmargin=1.6em,itemsep=2pt]
\item Ordnen Sie zu und lösen Sie: $x^{2} = 5x$; $x^{2} + 21 = 10x$; $x^{2} = 3x + 4$.
\item Welcher Typ hat zwei positive Lösungen? Welcher hat eine Lösung, die al-Chwarizmi nicht sieht?
\item Warum braucht man heute nur einen Typ? Welche zwei Erfindungen machen aus sechs Typen einen?
\item Al-Chwarizmi gibt für Typ (5) die Lösung $x = \frac b2\pm\sqrt{\frac{b^{2}}4 - c}$ und bemerkt: Ist $\frac{b^{2}}4 < c$, so ist die Aufgabe unmöglich. Was sagt man heute?
\end{enumerate}
\end{karte}

\stupser{gruen}{Das Quadrat $x^{2}$, an zwei Seiten je ein Rechteck $5x$ angelegt, wird durch ein Eckquadrat $25$ zum Quadrat $(x + 5)^{2} = 64$.}{22pt}
\begin{karte}{gruen}{7.20}{\typwerk}{\meta{1}{10 min}{Band 7, Kap.\ 6, Aufgabe 6.1}}
\kartentitel{Quadrat und zehn Wurzeln}
\worum{\enquote{Ein Quadrat und zehn Wurzeln sind gleich neununddreißig Dirhem}: die Musteraufgabe des Buches, das der Algebra ihren Namen gab. Dazu der geometrische Beweis --- quadratische Ergänzung als Figur.}
\methoden{Quadratische Ergänzung, Geometrische Algebra, Rezept entschlüsseln}
Lösen Sie $x^{2} + 10x = 39$ mit al-Chwarizmis Rezept und seinem geometrischen Beweis. Zeichnen Sie die Figur.
\end{karte}

\stupser{gruen}{$x^{3} + 4x = 8$: Mit $a = p^{2}$, $p = 2$, schneidet Chayyam die Parabel $x^{2} = py$ mit einem Kreis, dessen Durchmesser $\frac b{p^{2}} = 2$ auf der $x$-Achse liegt.}{22pt}
\begin{karte}{gruen}{7.21}{\typwerk}{\meta{3}{20 min}{Band 7, Kap.\ 6, Aufgabe 6.2}}
\kartentitel{Chayyams Kegelschnitte}
\worum{$500$ Jahre vor Cardano löste Omar Chayyam kubische Gleichungen --- geometrisch, als Schnitt von Parabel und Kreis. Er fand die Lösung als Strecke, nicht als Zahl, und wusste, dass ihm die Formel fehlte.}
\methoden{Kegelschnitte, Geometrische Lösung, Kubische Gleichung}
Führen Sie Chayyams Verfahren für $x^{3} + 4x = 8$ durch: Welche Parabel, welcher Kreis, und warum ist die $x$-Koordinate des Schnittpunkts die Lösung? Bestimmen Sie den Wert numerisch.
\end{karte}

%%%%%%%%%%% Kapitel 7
\stupser{tief}{$(2 + i)^{3} = 8 + 12i + 6i^{2} + i^{3} = 8 + 12i - 6 - i = 2 + 11i$. Also $\sqrt[3]{2 + \sqrt{-121}} = 2 + i$.}{22pt}
\begin{karte}{tief}{7.22}{\typtypisch}{\meta{3}{20 min}{Band 7, Kap.\ 7}}
\kartentitel{Bombelli und die Wurzel aus $-121$}
\worum{$x^{3} = 15x + 4$ hat die Lösung $4$. Cardanos Formel liefert $\sqrt[3]{2 + \sqrt{-121}} + \sqrt[3]{2 - \sqrt{-121}}$. Bombelli ($1572$) rechnete mit diesen \enquote{unmöglichen} Zahlen weiter, als wären sie welche --- und bekam $4$. So kamen die komplexen Zahlen in die Welt: nicht aus $x^{2} = -1$, sondern aus einer reellen Lösung.}
\methoden{Komplexe Zahlen, Cardanische Formel, Casus irreducibilis}
\begin{enumerate}[label=(\alph*),leftmargin=1.6em,itemsep=2pt]
\item Prüfen Sie, dass $x = 4$ die Gleichung $x^{3} = 15x + 4$ löst, und bestimmen Sie die beiden anderen Lösungen.
\item Wenden Sie Cardanos Formel ($x = \sqrt[3]{\frac q2 + \sqrt{\frac{q^{2}}4 - \frac{p^{3}}{27}}} + \sqrt[3]{\frac q2 - \sqrt{\dots}}$ für $x^{3} = px + q$) an. Was steht unter der Wurzel?
\item Zeigen Sie $(2 + \sqrt{-1})^{3} = 2 + \sqrt{-121}$ und vollenden Sie Bombellis Rechnung.
\item Warum heißt dieser Fall \emph{casus irreducibilis}, und was bewies erst Wantzel $1843$ darüber?
\end{enumerate}
\end{karte}

\stupser{gruen}{$F_n = F_{n-1} + F_{n-2}$: die Paare des Vormonats plus ein Neugeborenes je Paar, das schon vor zwei Monaten da war.}{22pt}
\begin{karte}{gruen}{7.23}{\typwerk}{\meta{1}{10 min}{Band 7, Kap.\ 7, Aufgabe 7.1}}
\kartentitel{Fibonaccis Kaninchen}
\worum{Die Aufgabe aus dem \emph{Liber abbaci} ($1202$), die der berühmtesten Folge der Mathematik den Namen gab --- und die Fibonacci selbst für eine Rechenübung hielt.}
\methoden{Lineare Rekursion, Modellannahmen, Fibonacci}
Jemand setzt ein Kaninchenpaar in einen Garten. Jedes Paar bringt vom zweiten Monat an in jedem Monat ein neues Paar hervor. Wie viele Paare sind es nach einem Jahr? Welche Modellannahmen stecken darin, und welchen Quotienten nähern die Zahlen?
\end{karte}

\stupser{gruen}{$x = u + v$ mit $3uv = -p$: Dann wird $x^{3} + px + q = 0$ zu $u^{3} + v^{3} = -q$, $u^{3}v^{3} = -\frac{p^{3}}{27}$ --- Vieta für $u^{3}, v^{3}$.}{22pt}
\begin{karte}{gruen}{7.24}{\typwerk}{\meta{2}{20 min}{Band 7, Kap.\ 7, Aufgabe 7.2}}
\kartentitel{Cardanos Formel}
\worum{Die erste Formel, die über die Antike hinausging: del Ferro fand sie, Tartaglia fand sie wieder, Cardano druckte sie ($1545$). Der Trick $x = u + v$ mit frei wählbarer Nebenbedingung ist bis heute Lehrstoff.}
\methoden{Cardanische Formel, Substitution (Gleichung), Vieta, Kubische Gleichung}
Leiten Sie die Cardanische Formel für $x^{3} + px + q = 0$ her und lösen Sie damit $x^{3} + 6x = 20$ (Cardanos eigenes Beispiel).
\end{karte}

\stupser{gruen}{Napier dachte kinematisch: ein Punkt bewegt sich mit abnehmender Geschwindigkeit, einer gleichmäßig. Sein Logarithmus ist $\mathrm{NapLog}(x) = 10^{7}\ln\frac{10^{7}}x$ --- kein $\log$ zur Basis $e$.}{22pt}
\begin{karte}{gruen}{7.25}{\typwerk}{\meta{2}{15 min}{Band 7, Kap.\ 7, Aufgabe 7.3}}
\kartentitel{Napiers Logarithmen}
\worum{$1614$: Multiplikation wird Addition. Napiers Tafeln halbierten die Rechenzeit der Astronomen --- Laplace sagte, sie hätten ihr Leben verdoppelt. Hier das Prinzip und eine Rechnung, wie sie Kepler machte.}
\methoden{Logarithmengesetze, Tabellenrechnen, Kinematische Definition}
Erklären Sie Napiers Grundprinzip und berechnen Sie mit vierstelligen Logarithmen (Basis $10$, Briggs) das Produkt $3{,}472\cdot2{,}718$ und die Wurzel $\sqrt[3]{1234}$. Wie genau ist das Ergebnis?
\end{karte}

%%%%%%%%%%% Kapitel 8
\stupser{tief}{Kreis um $(0, c)$ durch $(1, 1)$, geschnitten mit $y = x^{2}$: eine quadratische Gleichung in $y$. Berührung heißt doppelte Nullstelle --- das legt $c$ fest.}{22pt}
\begin{karte}{tief}{7.26}{\typtypisch}{\meta{3}{20 min}{Band 7, Kap.\ 8}}
\kartentitel{Descartes' Normalenmethode}
\worum{Descartes ($1637$) fand Tangenten ohne Grenzwert: Ein Kreis, der die Kurve berührt, hat mit ihr eine doppelte Nullstelle. Reine Algebra --- und der Grund, warum die \emph{Géométrie} keine Infinitesimalien braucht.}
\methoden{Doppelte Nullstelle, Tangente, Koordinatengeometrie}
Bestimmen Sie mit Descartes' Methode die Normale und die Tangente an die Parabel $y = x^{2}$ im Punkt $(1, 1)$: Kreis mit Mittelpunkt $(0, c)$ auf der Achse, der die Parabel in $(1, 1)$ berührt. Vergleichen Sie mit der Ableitung. Warum wird die Methode für andere Kurven schnell unhandlich?
\end{karte}

\stupser{gruen}{$x(b - x)\approx(x + e)(b - x - e)$; ausmultiplizieren, durch $e$ teilen, dann $e = 0$ setzen. Fermat sagt \enquote{adäqual}, nicht \enquote{gleich}.}{22pt}
\begin{karte}{gruen}{7.27}{\typwerk}{\meta{2}{10 min}{Band 7, Kap.\ 8, Aufgabe 8.1}}
\kartentitel{Fermats Adäqualität}
\worum{Um $1636$, dreißig Jahre vor Newton: Fermat findet Extremwerte, indem er $f(x)$ und $f(x + e)$ \enquote{adäqual} setzt, durch $e$ teilt und $e$ streicht. Das ist die Ableitung --- ohne Begriff dafür.}
\methoden{Extremwertaufgabe, Differenzenquotient, Infinitesimalien}
Teilen Sie mit Fermats Methode eine Strecke $b$ so, dass das Rechteck aus den Teilen maximal wird. Wo ist der Schritt, den Berkeley später \enquote{Geister verstorbener Größen} nannte?
\end{karte}

\stupser{gruen}{Fermat: alle Verläufe von zwei weiteren Runden aufzählen (auch die nicht gespielten) --- vier gleich wahrscheinliche Fälle. Pascal: vom Stand $2 : 2$ rückwärts.}{22pt}
\begin{karte}{gruen}{7.28}{\typwerk}{\meta{2}{15 min}{Band 7, Kap.\ 8, Aufgabe 8.2}}
\kartentitel{Das Teilungsproblem}
\worum{Zwei Spieler, je $32$ Pistolen, drei Siege gewinnen; Abbruch bei $2 : 1$. Pacioli teilte $2 : 1$, Tartaglia anders --- alle falsch. Pascal und Fermat lösten es $1654$ in Briefen, und die Wahrscheinlichkeitsrechnung begann.}
\methoden{Erwartungswert, Fallaufzählung, Rückwärtsrechnen}
Lösen Sie das Teilungsproblem mit Fermats und mit Pascals Methode. Warum ist Paciolis Antwort $2 : 1$ falsch --- was hat er gemessen?
\end{karte}

%%%%%%%%%%% Kapitel 9
\stupser{tief}{$\arctan x = x - \frac{x^{3}}3 + \frac{x^{5}}5 - \dots$ bei $x = 1$. Fehler nach $n$ Gliedern kleiner als das nächste Glied $\frac1{2n + 1}$.}{22pt}
\begin{karte}{tief}{7.29}{\typtypisch}{\meta{2}{15 min}{Band 7, Kap.\ 9}}
\kartentitel{Die Leibniz-Reihe}
\worum{$\frac\pi4 = 1 - \frac13 + \frac15 - \dots$ ($1674$): die schönste Reihe für $\pi$ und die unbrauchbarste. Machin ($1706$) nahm dieselbe Reihe an anderer Stelle und bekam $100$ Stellen.}
\methoden{Potenzreihe, Leibniz-Kriterium, Konvergenzgeschwindigkeit}
\begin{enumerate}[label=(\alph*),leftmargin=1.6em,itemsep=2pt]
\item Leiten Sie die Reihe aus $\arctan x = \int_0^x\frac{dt}{1 + t^{2}}$ her.
\item Wie viele Glieder braucht man für $\pi$ auf zwei Nachkommastellen? Auf zehn?
\item Machin: $\frac\pi4 = 4\arctan\frac15 - \arctan\frac1{239}$. Berechnen Sie $\pi$ damit auf sechs Stellen von Hand (je drei Glieder).
\item Warum konvergiert die Reihe bei $x = \frac15$ so viel schneller als bei $x = 1$?
\end{enumerate}
\end{karte}

\stupser{akzent}{$\frac1{1 + x} = 1 - x + x^{2} - \dots$ gilt für $|x| < 1$. Euler setzt $x = 1$ ein und bekommt $1 - 1 + 1 - \dots = \frac12$. Die Partialsummen sind $1, 0, 1, 0, \dots$}{22pt}
\begin{karte}{akzent}{7.30}{\typstolper}{\meta{2}{15 min}{Band 7, Kap.\ 9}}
\kartentitel{Was Euler durfte}
\worum{Euler rechnete mit divergenten Reihen, als wären sie konvergent, und bekam meist richtige Ergebnisse. Das 19.\ Jahrhundert verbot es; das 20.\ erlaubte es wieder --- unter Bedingungen. Die Aufgabe zeigt, was schiefgehen kann und was Euler intuitiv richtig machte.}
\methoden{Divergente Reihen, Gegenbeispiel, Abel-Summation}
\begin{enumerate}[label=(\alph*),leftmargin=1.6em,itemsep=2pt]
\item Euler: $1 - 1 + 1 - 1 + \dots = \frac12$ (aus $\frac1{1 + x}$ bei $x = 1$). Zeigen Sie, dass die Reihe divergiert, und dass man durch Klammern $0$ oder $1$ erhält.
\item Euler: $1 - 2 + 3 - 4 + \dots = \frac14$ (aus $\frac1{(1 + x)^{2}}$). Und $1 + 2 + 4 + 8 + \dots = -1$ (aus $\frac1{1 - x}$ bei $x = 2$). Was ist an diesen \enquote{Gleichungen} falsch, was ist richtig gemeint?
\item Abel-Summe: $\lim_{x\to1^{-}}\sum a_nx^{n}$. Berechnen Sie sie für $1 - 1 + 1 - \dots$ und zeigen Sie, dass sie für konvergente Reihen mit der Summe übereinstimmt (Abelscher Grenzwertsatz, ohne Beweis).
\item Formulieren Sie die Regel, die Euler fehlte: Wann darf man in eine Potenzreihe einen Randpunkt einsetzen?
\end{enumerate}
\end{karte}

\stupser{gruen}{$(1 + x)^{1/2} = 1 + \frac12x + \frac{\frac12\cdot(-\frac12)}{2!}x^{2} + \dots$; für $\sqrt2$ nimmt Newton nicht $x = 1$, sondern $\sqrt2 = \frac75\sqrt{\frac{50}{49}} = \frac75\sqrt{1 + \frac1{49}}$.}{22pt}
\begin{karte}{gruen}{7.31}{\typwerk}{\meta{2}{15 min}{Band 7, Kap.\ 9, Aufgabe 9.1}}
\kartentitel{Newtons Binomialreihe}
\worum{$1665$: Newton verallgemeinert den binomischen Satz auf gebrochene Exponenten --- ohne Beweis, durch Mustererkennung im Pascalschen Dreieck --- und berechnet damit Wurzeln und Logarithmen auf viele Stellen.}
\methoden{Binomialreihe, Potenzreihe, Geschickter Ansatz}
Leiten Sie mit Newtons verallgemeinertem binomischen Satz eine Reihe für $\sqrt{1 + x}$ her und berechnen Sie $\sqrt2$ auf sieben Stellen. Warum ist $x = 1$ eine schlechte und $x = \frac1{49}$ eine gute Wahl?
\end{karte}

\stupser{gruen}{$d(uv) = (u + du)(v + dv) - uv = u\,dv + v\,du + du\,dv$. Leibniz streicht $du\,dv$ als \enquote{unendlich klein von zweiter Ordnung}.}{22pt}
\begin{karte}{gruen}{7.32}{\typwerk}{\meta{1}{10 min}{Band 7, Kap.\ 9, Aufgabe 9.2}}
\kartentitel{Leibniz' Produktregel}
\worum{Die Produktregel, wie Leibniz sie $1675$ fand --- in einer Rechnung, die heute niemand mehr so aufschreiben dürfte und die trotzdem stimmt. Dazu die Frage, warum.}
\methoden{Produktregel, Infinitesimalien, Differentiale}
Leiten Sie mit Leibniz' Differentialkalkül die Produktregel her und beurteilen Sie das Argument: Wo ist es unstreng, wie rettet es der Grenzwert, wie rettet es die Nichtstandardanalysis?
\end{karte}

\stupser{gruen}{$\frac{\sin x}x = \prod\bigl(1 - \frac{x^{2}}{n^{2}\pi^{2}}\bigr)$; Koeffizient von $x^{2}$ links $-\frac16$, rechts $-\sum\frac1{n^{2}\pi^{2}}$.}{22pt}
\begin{karte}{gruen}{7.33}{\typwerk}{\meta{3}{20 min}{Band 7, Kap.\ 9, Aufgabe 9.3}}
\kartentitel{Das Basler Problem}
\worum{$\sum\frac1{n^{2}} = \frac{\pi^{2}}6$: Euler, $1735$, mit einem Argument, das er selbst für unsicher hielt --- er behandelte $\sin x$ wie ein Polynom. Hundert Jahre später wurde es durch Weierstraß' Produktsatz gerechtfertigt.}
\methoden{Produktdarstellung, Koeffizientenvergleich, Heuristik und Beweis}
Vollziehen Sie Eulers Lösung nach: $\sin x$ als unendliches Produkt über seine Nullstellen, Koeffizientenvergleich mit der Potenzreihe. Welche Lücke hat das Argument, und wie prüfte Euler sein Ergebnis numerisch?
\end{karte}

%%%%%%%%%%% Kapitel 10
\stupser{tief}{$A = \{r\in\mathbb Q\colon r\le0\text{ oder }r^{2} < 2\}$ hat kein größtes Element: Zu $r > 0$ mit $r^{2} < 2$ ist $r' = r + \frac{2 - r^{2}}{r + 2}$ größer und erfüllt $r'^{2} < 2$.}{22pt}
\begin{karte}{tief}{7.34}{\typtypisch}{\meta{2}{15 min}{Band 7, Kap.\ 10}}
\kartentitel{Ein Dedekind-Schnitt}
\worum{Dedekind ($1872$): Eine reelle Zahl ist ein Schnitt der rationalen Zahlen in Unter- und Oberklasse. $\sqrt2$ ist der Schnitt, bei dem die Lücke klafft --- und die Lücke selbst wird zur Zahl.}
\methoden{Dedekind-Schnitt, Vollständigkeit, Supremum und Infimum}
$A = \{r\in\mathbb Q\colon r\le0\text{ oder }r^{2} < 2\}$, $B = \mathbb Q\setminus A$.
\begin{enumerate}[label=(\alph*),leftmargin=1.6em,itemsep=2pt]
\item Zeigen Sie: $A$ ist nichtleer, nach oben beschränkt, und jedes Element von $A$ ist kleiner als jedes von $B$.
\item Zeigen Sie: $A$ hat kein größtes Element. (Zu $r\in A$, $r > 0$, finde ein größeres in $A$.)
\item Zeigen Sie: $B$ hat kein kleinstes Element. Warum ist genau das die \enquote{Lücke}?
\item Vergleichen Sie mit dem Schnitt $A' = \{r\colon r < 1\}$, $B' = \{r\colon r\ge1\}$. Was ist dort anders?
\end{enumerate}
\end{karte}

\stupser{gruen}{Stetig: $\sum a^{n}\cos(b^{n}\pi x)$ ist gleichmäßig konvergent (Majorante $\sum a^{n}$). Nicht differenzierbar: Die Ableitung der Summanden wächst wie $(ab)^{n}$ --- und $ab > 1$.}{22pt}
\begin{karte}{gruen}{7.35}{\typwerk}{\meta{2}{15 min}{Band 7, Kap.\ 10, Aufgabe 10.1}}
\kartentitel{Die Weierstraß-Funktion}
\worum{$1872$: eine überall stetige, nirgends differenzierbare Funktion. Hermite nannte sie eine \enquote{beklagenswerte Plage}. Sie zwang die Analysis, Stetigkeit und Differenzierbarkeit zu trennen.}
\methoden{Gleichmäßige Konvergenz, Gegenbeispiel, Majorante}
$W(x) = \sum_{n\ge0}a^{n}\cos(b^{n}\pi x)$ mit $0 < a < 1$, $b$ ungerade, $ab > 1 + \frac{3\pi}2$. Erklären Sie, warum $W$ stetig ist, und warum die Ableitung nirgends existiert (Idee, nicht der volle Beweis).
\end{karte}

\stupser{gruen}{Aus der Liste $x_1, x_2, \dots$ baut man eine Zahl, deren $n$-te Ziffer von der $n$-ten Ziffer von $x_n$ verschieden ist. Vorsicht bei $0{,}0999\ldots = 0{,}1$: Ziffern $0$ und $9$ vermeiden.}{22pt}
\begin{karte}{gruen}{7.36}{\typwerk}{\meta{2}{15 min}{Band 7, Kap.\ 10, Aufgabe 10.2}}
\kartentitel{Cantors Diagonale}
\worum{$1891$: ein Beweis auf einer Seite, dass es mehr reelle als natürliche Zahlen gibt. Das Argument ist so einfach, dass es misstrauisch macht --- und so stark, dass es Gödel und Turing trägt.}
\methoden{Diagonalargument, Widerspruchsbeweis, Abzählbarkeit}
Führen Sie Cantors Diagonalargument für $[0, 1]$ aus. Warum funktioniert dasselbe Argument nicht für $\mathbb Q$ --- wo bricht es? Und warum muss man bei Dezimalbrüchen mit $999\ldots$ aufpassen?
\end{karte}

%%%%%%%%%%% Kapitel 11
\stupser{tief}{$(x_1 - x_2)^{2}$ ändert sich nicht, wenn man $x_1$ und $x_2$ vertauscht --- also ist es durch die Koeffizienten ausdrückbar: $p^{2} - 4q$.}{22pt}
\begin{karte}{tief}{7.37}{\typtypisch}{\meta{2}{15 min}{Band 7, Kap.\ 11}}
\kartentitel{Lagranges Resolvente}
\worum{Lagrange ($1770$) fragte, warum die Formeln für Grad $2$, $3$, $4$ funktionieren --- und fand: Es geht um Ausdrücke in den Wurzeln, die unter Vertauschungen wenige Werte annehmen. Das ist der Keim der Galoistheorie.}
\methoden{Symmetrische Funktionen, Permutationen der Wurzeln, Vieta}
$x^{2} + px + q$ habe die Wurzeln $x_1, x_2$.
\begin{enumerate}[label=(\alph*),leftmargin=1.6em,itemsep=2pt]
\item Drücken Sie $x_1 + x_2$, $x_1x_2$ und $(x_1 - x_2)^{2}$ durch $p, q$ aus. Welche Rolle spielt die Symmetrie unter $x_1\leftrightarrow x_2$?
\item $x_1 - x_2$ ist nicht symmetrisch, aber $(x_1 - x_2)^{2}$ ist es. Folgern Sie die Lösungsformel.
\item Für $x^{3} + px + q$ mit Wurzeln $x_1, x_2, x_3$ und $\omega = e^{2\pi i/3}$ betrachte $t = x_1 + \omega x_2 + \omega^{2}x_3$. Wie viele verschiedene Werte nimmt $t^{3}$ unter den $6$ Permutationen an? Was folgt daraus?
\item Warum scheitert dieselbe Idee bei Grad $5$? (Nur die Zahl der Werte zählen; Band 11.)
\end{enumerate}
\end{karte}

\stupser{gruen}{Zerfällungskörper $\mathbb Q(\sqrt[3]2, \omega)$ hat Grad $6$; die Galoisgruppe permutiert die drei Wurzeln beliebig: $S_3$, auflösbar mit Normalteiler $A_3$.}{22pt}
\begin{karte}{gruen}{7.38}{\typwerk}{\meta{3}{20 min}{Band 7, Kap.\ 11, Aufgabe 11.1}}
\kartentitel{Die Galoisgruppe von $x^{3} - 2$}
\worum{Das kleinste Beispiel mit nichtabelscher Galoisgruppe --- und doch auflösbar, wie jede Gruppe der Ordnung $6$. Hier sieht man, was Galois' Begriff leistet und was Cardanos Formel damit zu tun hat.}
\methoden{Galoisgruppe, Zerfällungskörper, Auflösbarkeit}
Bestimmen Sie die Galoisgruppe von $x^{3} - 2$ über $\mathbb Q$ und prüfen Sie die Auflösbarkeit. Welche Zwischenkörper gehören zu den Untergruppen?
\end{karte}

%%%%%%%%%%% Kapitel 12
\stupser{tief}{Definition der Addition: $a + 0 = a$, $a + S(b) = S(a + b)$. Mit $2 = S(1)$, $1 = S(0)$: $2 + 2 = 2 + S(1) = S(2 + 1) = S(S(2 + 0)) = S(S(2))$.}{22pt}
\begin{karte}{tief}{7.39}{\typtypisch}{\meta{2}{15 min}{Band 7, Kap.\ 12}}
\kartentitel{$2 + 2 = 4$ nach Peano}
\worum{Peano ($1889$) reduziert die Arithmetik auf fünf Axiome und eine Nachfolgerfunktion. Dass $2 + 2 = 4$ dann ein \emph{Satz} ist, der bewiesen werden muss, zeigt, was Axiomatisierung bedeutet.}
\methoden{Peano-Axiome, Rekursive Definition, Vollständige Induktion}
\begin{enumerate}[label=(\alph*),leftmargin=1.6em,itemsep=2pt]
\item Definieren Sie $1, 2, 3, 4$ über die Nachfolgerfunktion $S$ und die Addition rekursiv. Beweisen Sie $2 + 2 = 4$.
\item Beweisen Sie $a + 1 = S(a)$ für alle $a$.
\item Beweisen Sie durch Induktion $0 + a = a$. Warum ist das nicht dasselbe wie das Axiom $a + 0 = a$?
\item Welches der fünf Axiome ist das Induktionsaxiom, und was passiert, wenn man es weglässt? (Nennen Sie ein Modell mit \enquote{zusätzlichen} Zahlen.)
\end{enumerate}
\end{karte}

%%%%%%%%%%% Kapitel 13
\stupser{akzent}{$R = \{x\colon x\notin x\}$. Ist $R\in R$? Dann $R\notin R$. Ist $R\notin R$? Dann $R\in R$. Beides ist unmöglich --- also gibt es $R$ nicht.}{22pt}
\begin{karte}{akzent}{7.40}{\typstolper}{\meta{2}{15 min}{Band 7, Kap.\ 13}}
\kartentitel{Die Menge aller Mengen}
\worum{Russell ($1901$) zeigte mit vier Zeilen, dass Cantors naive Mengenlehre widersprüchlich ist. Frege, dem er es schrieb, sah sein Lebenswerk zusammenbrechen. Die Lösung (Zermelo $1908$) ist eine Vorschrift, welche Mengen man bilden darf.}
\methoden{Russellsche Antinomie, Widerspruchsbeweis, Axiomatik}
\begin{enumerate}[label=(\alph*),leftmargin=1.6em,itemsep=2pt]
\item Führen Sie Russells Argument für $R = \{x\colon x\notin x\}$ aus.
\item Zeigen Sie mit Cantors Satz ($|M| < |\mathcal P(M)|$), dass es keine Menge aller Mengen gibt.
\item Zermelos Aussonderungsaxiom erlaubt nur $\{x\in M\colon\varphi(x)\}$ für eine schon vorhandene Menge $M$. Was wird dann aus $R$? Führen Sie das Argument für $R_M = \{x\in M\colon x\notin x\}$ durch.
\item Der Barbier, der alle rasiert, die sich nicht selbst rasieren: Was ist die Auflösung, und warum ist sie bei Mengen weniger bequem?
\end{enumerate}
\end{karte}
